%% ============================================================
%% AULA 14 — Encerramento: o mapa da inteligência computacional
%% GBC073 — Inteligência Computacional (FACOM/UFU)
%%
%% Esqueleto com esboço de conteúdo — a desenvolver.
%% Ementa (ficha SEI 5116656): fechamento do programa (itens 1 e 2
%% revistos). Os demais paradigmas (evolutiva, nebulosa,
%% bio-inspirados) ficam restritos à aula introdutória, por
%% diretriz do curso.
%%
%% Compilar:  latexmk -xelatex aula14.tex   (duas passadas!)
%% ============================================================
\documentclass[aspectratio=169, 11pt]{beamer}

\makeatletter
\def\input@path{{./}{../}}
\makeatother

\usetheme{InteligenciaComputacional}   % ANTES do babel: fontspec vem primeiro
\usepackage{babel}
\babelprovide[import, main]{brazilian}

\title[Aula 14 — Encerramento]{Aula 14: Encerramento — o mapa da inteligência computacional}
\subtitle[GBC073 — Inteligência Computacional]{Panorama \textbullet\ projeto final \textbullet\ próximos passos}
\author[Prof.\ Marcelo Albertini]{Prof.\ Marcelo Keese Albertini}
\institute{GBC073 — Inteligência Computacional \textbullet\ Universidade Federal de Uberlândia}
\date{2026}

\begin{document}

% ------------------------------------------------------------ capa
\begin{frame}[plain, noframenumbering]
  \titlepage
\end{frame}

% ------------------------------------------------------------ roteiro
\begin{frame}{Roteiro da aula}
  \begin{itemize}\setlength{\itemsep}{0.5ex}
    \item<1-> \textbf{Uma função, uma perda, um otimizador}
    \item<2-> \textbf{Do que vimos ao que o mercado usa}
    \item<3-> \textbf{Projeto final}
    \item<4-> \textbf{Próximos passos}
  \end{itemize}
\end{frame}

% ------------------------------------------------------------
\section{Uma função, uma perda, um otimizador}

\begin{frame}{Uma função, uma perda, um otimizador}
  \begin{itemize}\setlength{\itemsep}{0.4ex}
    \item O curso inteiro em uma frase: escolha uma família de funções,
          uma perda que mede o erro e um otimizador que desce o gradiente.
    \item Aula 1: a rede é um aproximador universal.
    \item Aulas 2--6: representar, aprender, otimizar, generalizar.
    \item Aulas 7--13: arquiteturas e aplicações — casos particulares do
          mesmo mapa.
  \end{itemize}
\end{frame}

% ------------------------------------------------------------
\section{Do que vimos ao que o mercado usa}

\begin{frame}{Do que vimos ao que o mercado usa}
  \begin{itemize}\setlength{\itemsep}{0.4ex}
    \item O que estudamos é a base do que os LLMs, sistemas de visão e
          recomendação usam.
    \item Transferência de aprendizado, ajuste fino, avaliação — o ciclo
          profissional inteiro.
    \item As ideias clássicas (perceptron, Hopfield, Kohonen) reaparecem
          dentro dos sistemas modernos.
  \end{itemize}
\end{frame}

% ------------------------------------------------------------
\section{Projeto final}

\begin{frame}{Projeto final}
  \begin{itemize}\setlength{\itemsep}{0.4ex}
    \item Escolher um problema real e aplicar o mapa: dados, modelo, perda,
          otimização, avaliação.
    \item Entregar código reproduzível e um relatório curto.
    \item Critérios: clareza da pergunta, honestidade da avaliação,
          reprodutibilidade.
  \end{itemize}
\end{frame}

% ------------------------------------------------------------
\section{Próximos passos}

\begin{frame}{Próximos passos}
  \begin{itemize}\setlength{\itemsep}{0.4ex}
    \item Disciplinas que continuam daqui: visão, NLP, aprendizado por
          reforço (GBC063).
    \item Para estudar sozinho: os cursos e livros citados ao longo do
          semestre.
    \item A regra de ouro: implemente do zero o que quiser entender —
          e use o pronto para o que quiser construir.
  \end{itemize}
\end{frame}

\end{document}
